\section{Conclusion}
\label{sec:conclusion}

In this work, we introduced the Business Process Scheduling Problem (BPSP) to address the challenges of scheduling business processes with stochastic activity durations and planned resource unavailability. Our solution framework combines process mining for parameter inference, deterministic transformations for uncertainty modeling, constraint programming for robust scheduling, and Monte Carlo simulation for probabilistic evaluation. Through evaluation on synthetic datasets, we demonstrated the adaptability of our solution to varying uncertainty levels, problem sizes, and resource configurations, achieving effective optimization of the probabilistic Makespan. Real-world validation on hospital data showcased the utility of the approach, achieving up to 14\% reductions in global Makespan.

Future work will focus on addressing the limitations identified in this study. Enhancing the scalability of the framework to handle larger, more complex business processes and integrating adaptive real-time rescheduling capabilities will extend its applicability to realistic environments. Exploring multi-objective optimization to balance trade-offs such as cost, resource utilization, and Makespan could further align the framework with organizational goals when scheduling business processes. Additionally, refining parameter inference methods to handle incomplete or noisy event logs and incorporating contextual data would significantly improve the practicality of the approach across diverse domains.
Lastly, while we currently use Monte Carlo simulation to handle uncertainty by sampling from the full distribution of activity durations, we aim to explore stochastic programming in future work. Stochastic programming offers a more structured approach by optimizing over a finite set of scenarios with assigned probabilities, potentially reducing uncertainty more effectively. However, its main limitation lies in the need to predefine a manageable set of scenarios, which is impractical in our case due to the vast number of possible realizations. As we introduce new uncertainties—such as variability in activity sequences—scenario-based methods may become more tractable, making stochastic programming a promising direction for future exploration.


\subsubsection{Acknowledgements}
We acknowledge the support of the PNRR project FAIR - Future AI Research (PE00000013), under the NRRP MUR program funded by the NextGenerationEU, the HORIZON 2020 HumanE-AI project (Grant 952026), and the PRIN project PINPOINT Prot. 2020FNEB27, CUP H23C22000280006 and H45E2100021000. 